\documentclass[11pt]{amsart}
\usepackage[T1]{fontenc}
\usepackage[utf8]{inputenc}
\usepackage[british]{babel}
\usepackage[margin=1in]{geometry}
\parindent=8pt
\parskip=3pt
\linespread{1.1}
\usepackage{amsmath,amssymb,amsthm,palatino,euler,microtype}
\renewcommand{\updefault}{n}
\usepackage{booktabs,longtable,array,enumitem,xcolor,upquote}
\usepackage[colorlinks=true,linkcolor=red!60!black,citecolor=blue!60!black,urlcolor=blue!60!black]{hyperref}
\hypersetup{pdftitle={Finite closures and graph certificates},pdfauthor={Espinoza, Frias-Armenta, Hernandez-Amador and Hernandez-Hernandez}}
\setlength{\emergencystretch}{2em}
\newcommand{\I}{\mathfrak I}
\newcommand{\IS}{\mathfrak I_S}
\newcommand{\ISV}{\mathfrak I_{SV}}
\newcommand{\ZZ}{\mathbb Z}
\newcommand{\cl}{\operatorname{cl}}
\newcommand{\Ga}{G_{14}^{a}}
\newcommand{\Gb}{G_{14}^{b}}
\newcommand{\Gc}{G_{14}^{c}}
\newcommand{\Dunce}{D_{15}}
\newcommand{\NS}{\mathcal N}
\newcommand{\BS}{\mathcal B}
\newtheorem{theorem}{Theorem}[section]
\newtheorem{proposition}[theorem]{Proposition}
\newtheorem{lemma}[theorem]{Lemma}
\newtheorem{corollary}[theorem]{Corollary}
\theoremstyle{definition}
\newtheorem{definition}[theorem]{Definition}
\newtheorem{remark}[theorem]{Remark}
\title[Finite closures and graph certificates]{Finite closures and graph certificates\\
Supplement to \emph{Counterexamples and enumeration for contractible and dismantlable graphs}}
\author[J. F. Espinoza]{Jes\'us F. Espinoza}
\author[M. E. Fr\'ias-Armenta]{Mart\'in Eduardo Fr\'ias-Armenta}
\author[R. G. Hern\'andez-Amador]{Rosal\'ia G. Hern\'andez-Amador}
\author[H. A. Hern\'andez-Hern\'andez]{H\'ector A. Hern\'andez-Hern\'andez}
\address{Departamento de Matem\'aticas, Universidad de Sonora, Hermosillo, Sonora, M\'exico}
\date{}
\begin{document}
\begin{abstract}
We describe the finite admissible-edge closures associated with the
counterexamples in the accompanying article. We prove the stopping rule
and the exhaustiveness of the closure algorithm, explain the natural
unfoldings of the 51 stuck seeds, and give an integral first-homology
presentation for their finite verification. We also record the
20 strictly 2-dismantlable graphs of order 12 and relate the earlier
15-vertex example to the present labelling.
\end{abstract}
\maketitle

\section{Graph classes and the finite constructions}
All graphs are finite and simple. The clique complex of a graph $G$ is
denoted by $\cl(G)$. We use induced open neighbourhoods $N_G(v)$ and
induced common neighbourhoods $N_G(u,v)$. All homology groups below have
integer coefficients unless a field is specified.

The class $\IS$ of strong I-contractible graphs contains $K_1$ and does
not contain the empty graph. For a graph of order at least two,
$G\in\IS$ precisely when either
\[
 N_G(v),\,G-v\in\IS\quad\hbox{for some }v,
 \qquad\hbox{or}\qquad
 N_G(u,v),\,G-uv\in\IS\quad\hbox{for some }uv\in E(G).
\]
The definition is well founded because each argument has smaller
$|V|+|E|$. A deletion, or its inverse addition, is \emph{admissible}
when the tested neighbourhood lies in $\IS$.
In the larger class $\I$, additions and deletions are both permitted,
and their neighbourhoods are required to lie in $\I$.
Starting from $K_1$ gives the least class closed under these four
operations \cite{DEFH2023}. Reversing a strong reduction shows
$\IS\subseteq\I$.

A graph of order at least two is \emph{stuck} if its clique complex is
acyclic over $\mathbb F_2$ and no vertex or edge is admissible for
deletion. The order restriction excludes $K_1$. A stuck graph is
outside $\IS$, since a strong reduction cannot start.

The accompanying article \cite{companion2026} classifies 51 stuck
$K_4$-free graphs of order 14. Three are I-contractible:
$\Ga,\Gb,\Gc$, with respectively 44, 45 and 45 edges. Each has an
admissible one-edge extension in $\IS$. The other 48 have
$H_1\cong\ZZ/3$ and no other reduced homology.
Their completeness as seeds depends on the enumeration and reduction
lemmas in the article. The finite constructions below start with
these explicitly deposited graphs. They permit edges that create a
$K_4$, unless a restriction is stated.

\section{Admissible-edge closures}
\begin{definition}
For a graph $G\notin\IS$, let $\NS(G)$ be the set of isomorphism classes
of graphs outside $\IS$ obtainable from $G$ by a finite, possibly empty,
sequence of admissible edge additions. Let $\BS(G)$ consist of the classes $[H+uv]\in\IS$
for which $[H]\in\NS(G)$ and $uv$ is admissible in $H$.
We call $\BS(G)$ the \emph{first-strong boundary}.
For a set of seeds $\mathcal S$, take the unions
$\NS(\mathcal S)$ and $\BS(\mathcal S)$.
\end{definition}
The boundary records the first entry into the strong class along some
path. A class may have several predecessors or paths from a seed.

\begin{lemma}\label{lem:monotone}
If $uv$ is a missing edge of $H$, then
$N_{H+uv}(u,v)=N_H(u,v)$ as induced graphs.
If $H\in\IS$ and this common neighbourhood belongs to $\IS$,
then $H+uv\in\IS$.
\end{lemma}
\begin{proof}
Neither endpoint belongs to the common neighbourhood. Adding $uv$
changes no adjacency within it and creates no new common neighbour.
Deleting $uv$ from $H+uv$ leaves $H$, so the edge alternative in the
definition of $\IS$ applies.
\end{proof}
\begin{corollary}\label{cor:no-return}
Every intermediate graph of an admissible-addition path ending outside
$\IS$ also lies outside $\IS$.
\end{corollary}
\begin{proof}
Once a path reaches $\IS$, Lemma~\ref{lem:monotone} applies to every
subsequent addition.
\end{proof}

The \emph{pruned closure algorithm} maintains a list of classes already
seen and a queue of non-strong classes still to be expanded. Initially
both the seen list and the queue contain the distinct non-strong seed
classes. On removing a class $[H]$,
test every missing edge $uv$ of a representative. If $N_H(u,v)\in\IS$,
form $H+uv$. Insert its class if unseen. Place a non-strong result in
the queue; record a strong result on the boundary and do not expand it.
The tests for $\IS$ use the full current graph, including every
previously added edge.

\begin{theorem}\label{thm:closure}
For a finite set of seeds, the pruned closure algorithm terminates.
Its non-strong output is exactly $\NS(\mathcal S)$, and its strong
output is exactly $\BS(\mathcal S)$.
\end{theorem}
\begin{proof}
On each branch the order stays fixed and the number of edges increases
strictly. A branch from an $n$-vertex graph has at most
$\binom n2-|E(G)|$ steps. There are finitely many graph classes at
each seed order.

Each inserted state is reached by tested admissible additions, so
every reported class has the asserted origin.
Conversely, take a path to a non-strong target.
By Corollary~\ref{cor:no-return}, every state on it is expanded unless
its class has already been expanded.
An isomorphism $\varphi:H\to H'$ sends $N_H(u,v)$ isomorphically to
$N_{H'}(\varphi(u),\varphi(v))$ and extends to
$H+uv\cong H'+\varphi(u)\varphi(v)$. Thus identifying isomorphic states
preserves every possible successor class. Induction along the path
shows that the target is found. Applying the same argument to the
last non-strong predecessor proves the assertion about the boundary.
\end{proof}

If only $K_4$-free targets are wanted, candidates containing a $K_4$
may also be discarded: subsequent additions cannot remove that clique.
The same argument proves exhaustiveness in this restricted universe.

\begin{remark}
For a 14-vertex graph, a common neighbourhood has at most 12 vertices.
The article's classification through order 13 implies that $\I$ and
$\IS$ agree on all such neighbourhoods; members of $\I$ are connected.
Consequently I-admissible and $\IS$-admissible edge additions coincide
at order 14. This observation inherits the finite classification's
computational dependence. Theorem~\ref{thm:closure} itself needs only
the definition of $\IS$ and holds at every order.
\end{remark}

\begin{lemma}\label{lem:homology}
An admissible edge addition induces isomorphisms on integral homology
in every degree. It preserves membership in $\I$ in both directions.
\end{lemma}
\begin{proof}
Let $C=N_H(u,v)$ and $K=\cl(C)$. The added simplices give
\[
 \cl(H+uv)=\cl(H)\cup(\Delta^{uv}*K),\qquad
 \cl(H)\cap(\Delta^{uv}*K)=(\partial\Delta^{uv})*K.
\]
Since $C\in\IS\subseteq\I$, the nonempty complex $K$ is contractible
\cite{EFH2022,DochtermannMatsushita2025}.
The intersection is the suspension of $K$. Both it and the added
piece are contractible.
The reduced Mayer--Vietoris sequence gives the homology isomorphisms
induced by inclusion. The addition and reverse deletion are
I-transformations, since their common neighbourhood is the same
member of $\I$.
\end{proof}

\section{The closure catalogs}\label{sec:catalogs}
Applying Theorem~\ref{thm:closure} to the three contractible seeds gives
\begin{center}
\begin{tabular}{lrrr}
\toprule
Seed & $|\NS(G)|$ & $|\BS(G)|$ & $K_4$-free members of $\NS(G)$\\
\midrule
$\Ga$ &16&32&4\\
$\Gb$ &10&16&3\\
$\Gc$ &5&8&1\\
\bottomrule
\end{tabular}
\end{center}
These finite computations use the proved recursion for $\IS$ and
exact isomorphism tests. The three searches are regenerated by the
separate finite verifier. Their non-strong union has 29 classes,
rather than 31, because classes from different seeds can coincide.
All 29 belong to $\I$ by Lemma~\ref{lem:homology}.
Their eight $K_4$-free members form the positive exceptional set in
row 14 of the article's table. The remaining 21 contain a $K_4$;
their simplicial collapsibility is not classified here.

Every class in $\BS(\Ga)$ supplies a counterexample of the type used
for Conjecture~15 of \cite{DEFH2023}: it has a representative $H+uv\in\IS$ whose
admissible edge $uv$ can be deleted to give $H\notin\IS$.
The count 32 describes this boundary, not every graph with such a
deletion.

\begin{proposition}[Computer-assisted finite catalog]\label{prop:2271}
For the set $\mathcal S$ of 51 deposited seeds,
$\NS(\mathcal S)$ has 2,271 isomorphism classes. Exactly 29 belong
to $\I$, and the other 2,242 have $H_1\cong\ZZ/3$.
Exactly 277 classes in this union are $K_4$-free: eight belong to
$\I$ and 269 have first homology $\ZZ/3$.
\end{proposition}
\begin{proof}
The complete pruned search on the 51 seeds, followed by exact
isomorphism identification, produces 2,271 non-strong classes.
For the torsion-seed part, exhaustiveness uses the complete-state
classifier and canonical labelling. A separate finite verifier
reconstructs the three positive-seed closures.
This finite search has been regenerated with the current complete-state
classifier. Its class set agrees with the deposited catalog.
The stored paths are checked independently, as are nonmembership in
$\IS$, pairwise nonisomorphism and integral first homology.

The three contractible seeds contribute precisely the 29-class union
above. Every other seed has $H_1\cong\ZZ/3$.
Lemma~\ref{lem:homology} preserves that group along each addition path,
so none of those descendants belongs to $\I$.
The two parts are disjoint, with sizes $29$ and $2271-29=2242$.
Testing every four-vertex subset gives the restricted counts 8 and 269.
\end{proof}

The $K_4$-free non-strong exceptional set in row 14 has 633 classes.
Exact comparison places all 277 restricted closure classes inside
that set, leaving 356 classes outside these closures. One is encoded
by \texttt{M?B@\char96 jIja\{J\}jy\{]?}.
Thus the 2,271-class union does not exhaust even the $K_4$-free
exceptions at order 14.
The positive part is complete in the $K_4$-free universe by the
article's reduction theorem; the 29-class positive union is not claimed
to classify the unrestricted order-14 universe.

\section{Natural unfoldings of the stuck seeds}\label{sec:unfold}
For a $K_4$-free graph, the clique complex has dimension at most two.
Let $d(e)$ be the number of triangles containing an edge $e$.
Begin with a separate closed triangle for each triangle of $\cl(G)$.
Whenever $d(e)=2$, glue the two copies of $e$ preserving the endpoint
labels. Make no other identifications. The resulting complex is the
\emph{natural unfolding} $U(G)$.
When every edge and vertex belongs to a triangle, as for all 51 seeds,
the labels define a quotient map $U(G)\to\cl(G)$.
The construction identifies triangle corners as well as edges.
Distinct edges of the unfolding may have the same image labels.

\begin{proposition}\label{prop:unfold}
The triangle list determines the unfolding's components, vertex links
and boundary words by finite gluing.
If every vertex link is a circle or an interval, the unfolding is a
compact surface with boundary. A connected component with Euler
characteristic 1 and one boundary component is a disc; one with Euler
characteristic 0 and two boundary components is an annulus.
\end{proposition}
\begin{proof}
Represent a corner by $(T,v)$, where $T$ is a triangle containing $v$.
For an edge $uv$ in exactly two triangles $T,T'$, identify
$(T,u)$ with $(T',u)$ and $(T,v)$ with $(T',v)$.
The resulting equivalence classes are the quotient vertices.
The glued triangle edges are its edges. Joining triangle pairs along
the prescribed edges determines the connected components; no other
identification joins two components.

Each incident triangle contributes an interval to a vertex link.
Glue interval endpoints exactly when the corresponding triangle
edges were glued. A circular link gives a disc neighbourhood; an
interval link gives a half-disc neighbourhood. This proves the local
surface condition. Unglued triangle edges form the boundary;
traversing its cycles gives the boundary words in original labels.
Cell counts give the Euler characteristic.
The classification of compact connected surfaces now gives the
two stated types.
\end{proof}

All 51 unfoldings satisfy this surface condition. The cross-tabulation
with the homology of the original clique complex is
\begin{center}
\begin{tabular}{lrrr}
\toprule
Integral $H_1(\cl(G))$ &Disc&Annulus&Disc $\sqcup$ annulus\\
\midrule
$0$&3&0&0\\
$\ZZ/3$&45&2&1\\
\bottomrule
\end{tabular}
\end{center}
Thus the three contractible seeds have disc unfoldings, and there
are 45 further discs with torsion after boundary identification.
The 48/2/1 split is reconstructed from the corner gluing,
independently of the homology column.
The annuli have boundary lengths 4 and 12. In the disconnected
unfolding, the annulus has boundary lengths 8 and 12 and the disc has
boundary length 4. Table~\ref{tab:seeds} gives all encodings and
boundary lengths; in row 51 the length 4 belongs to the disc and
lengths 8 and 12 to the annulus. The data also record every cyclic boundary word
and vertex-link certificate.

\section{Integral first homology from triangle relations}
The finite checks use an integer presentation that applies to every
graph. Its output concerns first homology; higher homology requires
additional information.

\begin{proposition}\label{prop:h1}
Orient all edges of $G$ from smaller to larger endpoints, and let $F$
be a spanning forest of $G$. Put $Q=E(G)\setminus E(F)$.
Form an integer matrix $M$ with rows indexed by $Q$ and columns by
triangles. The column for $i<j<k$ records the coefficients of
$[j,k]-[i,k]+[i,j]$ on the edges in $Q$.
Then $H_1(\cl(G);\ZZ)\cong\operatorname{coker}M$.
If the nonzero Smith factors are $d_1,\ldots,d_r$, this group is
\[
 \ZZ^{\,|Q|-r}\oplus
 \bigoplus_{\substack{1\leq i\leq r\\d_i>1}}\ZZ/d_i.
\]
\end{proposition}
\begin{proof}
An edge outside $F$ and the unique forest path between its endpoints
determine a fundamental cycle. These cycles form an integral basis
of $\ker\partial_1$. Equivalently, recording coefficients on $Q$
gives an isomorphism $\ker\partial_1\to\ZZ^Q$: the remaining
coefficients are uniquely determined by the cycle equations on
each tree. Isolated vertices contribute no generator.
The boundary of each triangle maps to its column of $M$.
First homology is the cycle lattice modulo triangle boundaries.
Higher-dimensional simplices do not change this quotient.
The Smith normal form gives the displayed decomposition.
\end{proof}

The Smith computation uses unimodular integer row and column
operations and freezes a pivot only after it divides the remaining
block. For a connected two-dimensional complex with $\chi=1$ and
finite $H_1$, the group $H_2$ is free of rank zero by the Euler
characteristic. This gives the higher-homology conclusion for the
48 torsion seeds. For their descendants,
Lemma~\ref{lem:homology} supplies the conclusion in every dimension,
including when added edges create higher cliques.

\section{The strictly 2-dismantlable catalog}
A nonempty graph is a cone if a vertex is adjacent to every other
vertex. Cones form level $(-1)$-Dis. For $k\geq0$, a vertex is
$k$-dismantlable if its open neighbourhood belongs to $(k-1)$-Dis;
a graph belongs to $k$-Dis if successive such deletions reduce it
to $K_1$. The empty graph belongs to no level.
The hierarchy satisfies
$1\text{-Dis}\subseteq2\text{-Dis}\subseteq\ISV$,
where $\ISV$ is the strong class using vertex gluings alone
\cite{DEFH2023}.

\begin{proposition}[Computer-assisted]
The minimum order of a graph in $2$-Dis outside $1$-Dis is 12.
There are exactly 20 isomorphism classes of that order.
\end{proposition}
\begin{proof}
Vertex deletion cannot merge components, and an isolated vertex
cannot be deleted because its neighbourhood is empty. Thus every $2$-dismantlable graph is
connected, so the connected universes in the table suffice.
For orders 3 through 11, the article's table gives equal counts for
$1$-Dis and $\ISV$. The inclusions force equality with $2$-Dis.
At orders one and two, its only members are $K_1$ and $K_2$,
both in $1$-Dis. At order 12 the difference is
\[
 22\,666\,525\,713-22\,666\,525\,693=20.
\]
The 20 deposited graphs are pairwise nonisomorphic. Exact recursion
places each in $2$-Dis and outside $1$-Dis. They exhaust the difference
and establish existence at the minimum order.
Completeness and minimality inherit the table's enumeration dependence.
\end{proof}

Table~\ref{tab:strict} records the graphs and elementary invariants.
Automorphism-group orders are exact counts of adjacency-preserving
permutations, included as catalog invariants.

\section{The earlier 15-vertex example}
Let $\Dunce$ have encoding
\texttt{N\textbar thHsL\char96 sHoI\char96 RPOwLG}.
It has 15 vertices, 46 edges and 32 triangles.
The correspondence with Figure~2A of \cite{DEFH2023} is
\[
\begin{array}{c|rrrrrrrrrrrrrrr}
 \text{present}&0&1&2&3&4&5&6&7&8&9&10&11&12&13&14\\
 \text{published}&1&4&13&12&11&3&14&15&10&5&2&6&7&8&9
\end{array}
\]
Under this bijection the complete triangle lists agree. The image of
the full edge set also agrees with the union of edges in the recorded
published triangles. These comparisons identify the recorded
triangulation and its graph in the two labellings.
The missing edge $\{9,12\}$ corresponds to the missing edge
$\{5,7\}$ in the published labelling.

The unfolding is a disc with face vector $(23,54,32)$ and boundary word
\[
 (0,1,5,10,\;0,1,5,10,\;0,10,5,1).
\]
Its image is a four-cycle, and the attaching degree is one.
Collapsing a spanning tree of that cycle gives a circle with a
two-cell attached by degree one, so $\cl(\Dunce)$ is contractible.
The disc verifier checks interior injectivity and that all
identifications occur on the boundary.

For a direct graph certificate, add $\{9,12\}$.
The common neighbourhood is strong, and the resulting graph has
the strong reduction in Table~\ref{tab:d15}.
Every displayed neighbourhood is taken in the current graph;
its stored vertex reduction is checked recursively.
Conversely, the vertex links of $\Dunce$ are triangle-free and not
trees, and its edge links are edgeless graphs of order at least two.
None belongs to $\IS$, so $\Dunce$ is stuck.
This relates the earlier example to the present constructions,
without asserting a new minimum at order 15.

\section{Data and verification}
The encodings, triangle lists, boundary words and reductions are
in the companion repository:
\begin{center}\small
\url{https://github.com/Geometric-and-Combinatorial-Structures/minimal-I-contractible-not-strong}.
\end{center}
The article contains the essential counterexample proofs and the
algorithms for the table. This supplement supplies the mathematical
basis and scope of the auxiliary catalogs. Execution instructions,
requirements and provenance are given separately in the repository.

The finite verification reconstructs five disc presentations, all
51 unfoldings, the three small closures and the displayed reductions.
It checks the 2,271 stored paths and integral first homology groups,
and compares the catalogs up to isomorphism. The full 51-seed
closure has also been regenerated with the current complete-state
C classifier. This search on fixed seeds is distinct from the much
larger enumerations establishing the table and seed completeness.
Those retain their documented historical computational dependence;
the finite checks do not supply a full independent recount.

\appendix
\section{Catalog tables}
The machine-readable catalogs are the authoritative source for copying
graph6 encodings; PDF text extraction may substitute glyphs.
\begingroup\small
\setlength{\tabcolsep}{3pt}
\begin{longtable}{rlrrrrl}
\caption{The 51 stuck $K_4$-free seeds of order 14. Here $m$ counts edges, $t$ triangles, and $a=|\operatorname{Aut}(G)|$. The final column gives all boundary lengths of the natural unfolding. Rows 48 and 49 are annuli; row 51 is a disc together with an annulus; the other rows are discs.}\label{tab:seeds}\\
\toprule
No. & Encoding & $m$ & $t$ & $a$ & $H_1$ & Boundary \\
\midrule\endfirsthead
\toprule
No. & Encoding & $m$ & $t$ & $a$ & $H_1$ & Boundary \\
\midrule\endhead
\midrule\multicolumn{7}{r}{Continued on the next page}\\
\endfoot\bottomrule\endlastfoot
1 & \verb+M?b@bFHNEk@~|k}U?+ & 44 & 31 & 1 & $0$ & 15 \\
2 & \verb+M?B@thku?~Lylk}[?+ & 45 & 32 & 2 & $0$ & 16 \\
3 & \verb+M?r@d`ifbiFfo}~o?+ & 45 & 32 & 2 & $0$ & 18 \\
4 & \verb+M?B@cXamAu\r{]V{?+ & 43 & 30 & 2 & $\ZZ/3$ & 12 \\
5 & \verb+M?B@cXemAuB}ze{]?+ & 43 & 30 & 2 & $\ZZ/3$ & 12 \\
6 & \verb+M?BDA_wqpmJvj]|w?+ & 43 & 30 & 2 & $\ZZ/3$ & 12 \\
7 & \verb+M?`@E_wDf[UnU}z{?+ & 43 & 30 & 6 & $\ZZ/3$ & 12 \\
8 & \verb+M?`adQWT\LK^du^w?+ & 43 & 30 & 2 & $\ZZ/3$ & 12 \\
9 & \verb+M?`adQWtVDUr[]^w?+ & 43 & 30 & 2 & $\ZZ/3$ & 12 \\
10 & \verb+M?bB@_WmbRJZ|M~w?+ & 43 & 30 & 1 & $\ZZ/3$ & 12 \\
11 & \verb+M?bB@_WsbX^Jlm~w?+ & 43 & 30 & 1 & $\ZZ/3$ & 12 \\
12 & \verb+M?bB@_WsequuZZ~w?+ & 43 & 30 & 1 & $\ZZ/3$ & 12 \\
13 & \verb+M?bB@_Wu`ZJjlm~w?+ & 43 & 30 & 2 & $\ZZ/3$ & 12 \\
14 & \verb+M?bB@_WubXNJlm~w?+ & 43 & 30 & 2 & $\ZZ/3$ & 12 \\
15 & \verb+M?bB@aWYTkZZ^M|w?+ & 43 & 30 & 1 & $\ZZ/3$ & 12 \\
16 & \verb+M?bB@aWZSmXZ|wNm?+ & 43 & 30 & 1 & $\ZZ/3$ & 12 \\
17 & \verb+M?bB@aWZTkFf|wuu?+ & 43 & 30 & 1 & $\ZZ/3$ & 12 \\
18 & \verb+M?bB@aWdbYFv|wuu?+ & 43 & 30 & 1 & $\ZZ/3$ & 12 \\
19 & \verb+M?bB@aWdbZEv|wuu?+ & 43 & 30 & 1 & $\ZZ/3$ & 12 \\
20 & \verb+M?bB@aWqplQzZ]|w?+ & 43 & 30 & 1 & $\ZZ/3$ & 12 \\
21 & \verb+M?bB@aWyTkYZ^M|w?+ & 43 & 30 & 1 & $\ZZ/3$ & 12 \\
22 & \verb+M?bB@aXdbY[{Nmuu?+ & 43 & 30 & 1 & $\ZZ/3$ & 12 \\
23 & \verb+M?bB@aXlbYFfxwuu?+ & 43 & 30 & 2 & $\ZZ/3$ & 12 \\
24 & \verb+M?bB@aXyTkYZxw^M?+ & 43 & 30 & 1 & $\ZZ/3$ & 12 \\
25 & \verb+M?otU`SmVEKm]M^o?+ & 43 & 30 & 4 & $\ZZ/3$ & 12 \\
26 & \verb+M?qa``{rfEEymitY?+ & 43 & 30 & 4 & $\ZZ/3$ & 12 \\
27 & \verb+M?qa``{ufEUimitY?+ & 43 & 30 & 4 & $\ZZ/3$ & 12 \\
28 & \verb+M?r@daiRbrTYd{tb?+ & 43 & 30 & 1 & $\ZZ/3$ & 12 \\
29 & \verb+M?`Db_kVcyl]hyvw?+ & 44 & 31 & 2 & $\ZZ/3$ & 15 \\
30 & \verb+M?`Db_khi{L\nUvw?+ & 44 & 31 & 1 & $\ZZ/3$ & 15 \\
31 & \verb+M?`adL[mBkvLY{d{?+ & 44 & 31 & 2 & $\ZZ/3$ & 15 \\
32 & \verb+M?`cu_whRRFZ~Wtm?+ & 44 & 31 & 1 & $\ZZ/3$ & 15 \\
33 & \verb+M?`cv@UNDNZMtsY{?+ & 44 & 31 & 1 & $\ZZ/3$ & 15 \\
34 & \verb+M?b@bCw`usLvvw{]?+ & 44 & 31 & 1 & $\ZZ/3$ & 15 \\
35 & \verb+M?qa``[NNEYlme~o?+ & 44 & 31 & 2 & $\ZZ/3$ & 15 \\
36 & \verb+M?qa``[iUEylNm~o?+ & 44 & 31 & 1 & $\ZZ/3$ & 15 \\
37 & \verb+M?qa``[icjk^Nk~o?+ & 44 & 31 & 2 & $\ZZ/3$ & 15 \\
38 & \verb+M?qa``[iej[ZNk~o?+ & 44 & 31 & 2 & $\ZZ/3$ & 15 \\
39 & \verb+M?qa``kLeE^wt]i}?+ & 44 & 31 & 2 & $\ZZ/3$ & 15 \\
40 & \verb+M?qa``kMmMIniu~o?+ & 44 & 31 & 2 & $\ZZ/3$ & 15 \\
41 & \verb+M?qa``kifEfUtZ~o?+ & 44 & 31 & 1 & $\ZZ/3$ & 15 \\
42 & \verb+M?qa`awkmiNdwu~o?+ & 44 & 31 & 2 & $\ZZ/3$ & 15 \\
43 & \verb+M?qa`awwRxYj{k~o?+ & 44 & 31 & 2 & $\ZZ/3$ & 15 \\
44 & \verb+M?qa`g{guLzwmkU]?+ & 44 & 31 & 2 & $\ZZ/3$ & 15 \\
45 & \verb+M?qa`jSwubncvomd?+ & 44 & 31 & 2 & $\ZZ/3$ & 15 \\
46 & \verb+M?qaejaLe[ngmh{[?+ & 44 & 31 & 2 & $\ZZ/3$ & 15 \\
47 & \verb+M?`FSzcVBmVFs{J{?+ & 45 & 32 & 4 & $\ZZ/3$ & 18 \\
48 & \verb+M?qa``{ws}Lxmkt[?+ & 45 & 32 & 8 & $\ZZ/3$ & 4,12 \\
49 & \verb+M?qa``{wvkVUt[}Q?+ & 45 & 32 & 8 & $\ZZ/3$ & 4,12 \\
50 & \verb+MCOfczgVB]UeM{p{?+ & 45 & 32 & 4 & $\ZZ/3$ & 18 \\
51 & \verb+MCp`faiFf{W}][js?+ & 47 & 34 & 2 & $\ZZ/3$ & 4,8,12 \\
\end{longtable}\endgroup

\begingroup\small
\setlength{\tabcolsep}{3pt}
\begin{longtable}{rlrrrrrrr}
\caption{The 20 graphs of order 12 in $2$-Dis outside $1$-Dis. The columns give edges $m$, triangles $t$, four-cliques $q$, clique number $\omega$, degrees $\delta,\Delta$, and automorphism-group order $a$.}\label{tab:strict}\\
\toprule
No. & Encoding & $m$ & $t$ & $q$ & $\omega$ & $\delta$ & $\Delta$ & $a$ \\
\midrule\endfirsthead
\toprule
No. & Encoding & $m$ & $t$ & $q$ & $\omega$ & $\delta$ & $\Delta$ & $a$ \\
\midrule\endhead
\midrule\multicolumn{9}{r}{Continued on the next page}\\
\endfoot\bottomrule\endlastfoot
1 & \verb+KEhutxmv|~^{+ & 46 & 69 & 34 & 4 & 6 & 9 & 4 \\
2 & \verb+KEjer^{z}~^{+ & 47 & 73 & 37 & 4 & 6 & 9 & 2 \\
3 & \verb+KQyurzU~t~^{+ & 47 & 72 & 37 & 5 & 7 & 9 & 4 \\
4 & \verb+KQyurzU~u~^{+ & 47 & 73 & 39 & 5 & 6 & 9 & 2 \\
5 & \verb+KEhvtzm||~^{+ & 48 & 78 & 45 & 5 & 6 & 9 & 2 \\
6 & \verb+KEjeu~m^~z^r+ & 48 & 77 & 42 & 5 & 6 & 9 & 1 \\
7 & \verb+KEnfbzm||~^{+ & 48 & 75 & 38 & 4 & 7 & 9 & 4 \\
8 & \verb+KEnfbzm~]~^{+ & 48 & 75 & 38 & 4 & 7 & 9 & 2 \\
9 & \verb+KEnfbzm~^v^{+ & 48 & 75 & 38 & 4 & 7 & 9 & 2 \\
10 & \verb+KQyurzuv|~^{+ & 48 & 77 & 44 & 5 & 7 & 9 & 4 \\
11 & \verb+KQyurzuv~}^{+ & 48 & 76 & 42 & 5 & 7 & 9 & 2 \\
12 & \verb+KEht~jzz~}^f+ & 49 & 82 & 50 & 5 & 7 & 9 & 2 \\
13 & \verb+KEjflzz^~}^f+ & 49 & 81 & 47 & 5 & 7 & 9 & 2 \\
14 & \verb+KEnb~jl~t~^{+ & 49 & 81 & 48 & 5 & 7 & 9 & 2 \\
15 & \verb+KEnb~jl~v|^{+ & 49 & 80 & 45 & 5 & 8 & 9 & 2 \\
16 & \verb+KEnff^{^~v^{+ & 49 & 80 & 44 & 5 & 7 & 9 & 1 \\
17 & \verb+KEnfnZ{~r~^{+ & 49 & 79 & 42 & 5 & 8 & 9 & 4 \\
18 & \verb+KUZurzm~vv^{+ & 49 & 81 & 48 & 5 & 8 & 9 & 4 \\
19 & \verb+KEnb~rvz~}^f+ & 50 & 86 & 54 & 5 & 8 & 9 & 2 \\
20 & \verb+KEzftzv^~}^f+ & 50 & 85 & 51 & 5 & 8 & 9 & 2 \\
\end{longtable}\endgroup

\begingroup\small
\setlength{\tabcolsep}{3pt}
\begin{longtable}{rlll}
\caption{A strong reduction of $\Dunce+\{9,12\}$. The symbols $v$ and $e$ indicate vertex and edge deletions. The last column lists vertex deletions in the tested neighbourhood, followed by its surviving vertex in brackets. The main reduction ends at vertex $14$.}\label{tab:d15}\\
\toprule
Step & Delete & Current neighbourhood & Link reduction \\
\midrule\endfirsthead
\toprule
Step & Delete & Current neighbourhood & Link reduction \\
\midrule\endhead
\midrule\multicolumn{4}{r}{Continued on the next page}\\
\endfoot\bottomrule\endlastfoot
1 & $v_{\{11\}}$ & $\{0,7,9,12\}$ & $0,7,9\ ;\ [12]$ \\
2 & $e_{\{0,9\}}$ & $\{10,12\}$ & $10\ ;\ [12]$ \\
3 & $e_{\{0,12\}}$ & $\{10\}$ & $\ ;\ [10]$ \\
4 & $v_{\{10\}}$ & $\{0,5,8,9,12,14\}$ & $0,9,12,5,8\ ;\ [14]$ \\
5 & $v_{\{8\}}$ & $\{4,5,7,14\}$ & $5,4,7\ ;\ [14]$ \\
6 & $v_{\{4\}}$ & $\{0,1,3,5,7\}$ & $5,1,0,3\ ;\ [7]$ \\
7 & $v_{\{0\}}$ & $\{1,2,3,14\}$ & $3,2,1\ ;\ [14]$ \\
8 & $v_{\{1\}}$ & $\{2,5,13,14\}$ & $2,5,13\ ;\ [14]$ \\
9 & $v_{\{2\}}$ & $\{3,5,6,7\}$ & $3,5,6\ ;\ [7]$ \\
10 & $v_{\{3\}}$ & $\{7\}$ & $\ ;\ [7]$ \\
11 & $v_{\{5\}}$ & $\{6,9,12,13\}$ & $6,9,12\ ;\ [13]$ \\
12 & $v_{\{6\}}$ & $\{7,9\}$ & $7\ ;\ [9]$ \\
13 & $v_{\{7\}}$ & $\{9,12,13,14\}$ & $9,12,13\ ;\ [14]$ \\
14 & $v_{\{9\}}$ & $\{12\}$ & $\ ;\ [12]$ \\
15 & $v_{\{12\}}$ & $\{13\}$ & $\ ;\ [13]$ \\
16 & $v_{\{13\}}$ & $\{14\}$ & $\ ;\ [14]$ \\
\end{longtable}\endgroup

\bibliographystyle{siam}
\bibliography{refs}
\end{document}
